\section{总结与延伸}

本讲从 Okta 创业与运营经历提炼出 identity infrastructure 的可迁移原则：

\begin{enumerate}
  \item \textbf{Transition 需要 wedge}：技术变迁提供窗口，客户今天愿意付费的窄入口负责启动，平台深度负责长期价值；
  \item \textbf{Identity 是 control plane}：directory、authentication、authorization、provisioning、session 与 audit 共同决定访问；
  \item \textbf{Front door 同时要求安全和可用}：availability、latency、correctness、blast radius 与 safe change 必须联合验收；
  \item \textbf{信念必须连接 evidence}：低概率环境需要 daily progress、buyer signal 和诚实 uncertainty，而非拒绝数据；
  \item \textbf{治理依赖无失真信息}：CEO 与 board 要把坏消息转成 options、decision rights 和 follow-through；
  \item \textbf{Incident 必须关闭学习闭环}：containment、RCA、remediation、notification 和 culture change 缺一不可；
  \item \textbf{Zero Trust 建立可撤销信任}：上下文、least privilege、step-up 和 continuous re-evaluation 依赖干净 lifecycle；
  \item \textbf{Agent 使用 delegation}：独立 principal、actor/subject、short-lived scoped token、vault 与 audit 替代共享秘密；
  \item \textbf{并购先保护 customer contract}：共享 primitive 不代表立即合并 runtime、团队和产品工作流；
  \item \textbf{AI 从 demo 走向 workflow}：category value 需要能力、流程、组织和安全控制共同变化。
\end{enumerate}

\begin{importantbox}{Identity Architecture 作业}
为一个“AI agent 代表员工读取工单并创建代码草稿”的系统提交一页设计：画出 human、agent principal、authorization server、ticket system、code host 和 audit；定义 OIDC/OAuth/token-exchange flow、scope、audience、expiry、vault、session revoke 与 owner offboarding；再写出 IdP outage、stale group、token theft 和 agent compromise 的降级策略。禁止让 agent 直接继承用户长期 token。
\end{importantbox}

\subsection{拓展阅读}

{\small
\begin{itemize}[nosep,leftmargin=*]
  \item \href{https://developer.okta.com/docs/concepts/universal-directory/}{\textbf{Okta Universal Directory}}：统一 identity data model 与 source integration。
  \item \href{https://help.okta.com/oie/en-us/content/topics/identity-engine/policies/about-policies.htm}{\textbf{Okta Identity Engine policies}}：global session 与 app sign-in assurance。
  \item \href{https://sec.okta.com/articles/2023/11/unauthorized-access-oktas-support-case-management-system-root-cause/}{\textbf{Okta support-system incident RCA}}：timeline、logging gap 与 remediation。
  \item \href{https://csrc.nist.gov/pubs/sp/800/207/final}{\textbf{NIST SP 800-207}}：Zero Trust Architecture。
  \item \href{https://www.rfc-editor.org/rfc/rfc8693}{\textbf{RFC 8693}}：OAuth 2.0 Token Exchange。
  \item \href{https://www.rfc-editor.org/rfc/rfc7644}{\textbf{RFC 7644}}：SCIM provisioning protocol。
  \item \href{https://www.w3.org/TR/webauthn-3/}{\textbf{WebAuthn Level 3}}：phishing-resistant public-key authentication。
\end{itemize}
}

\end{document}
